\documentclass{article}
\usepackage[utf8]{inputenc}
\usepackage[T1]{fontenc}
\usepackage{geometry}
\geometry{
    a4paper,
    left=2.5cm,
    right=2.5cm,
    top=2.5cm,
    bottom=2.5cm
}
\usepackage{graphicx}
\usepackage{float}
\usepackage{hyperref}
\usepackage[english]{babel}
\usepackage{amsmath}
\usepackage{amssymb}
\usepackage{array}
\usepackage{booktabs}
\usepackage{caption}
\usepackage{subcaption}
\usepackage{enumitem}
\usepackage{fancyhdr}
\usepackage{titlesec}
\usepackage{url}
\usepackage{textcomp}
\usepackage{xcolor}
\usepackage{listings}
\usepackage{multirow}
\usepackage{longtable}
\usepackage{tabularx}
\usepackage{times}

% Set up page style
\pagestyle{fancy}
\fancyhf{}
\rfoot{\thepage}
\renewcommand{\headrulewidth}{0pt}

% Title formatting
\titleformat{\chapter}[display]{\normalfont\bfseries}{\chaptertitle\thechapter}{0pt}{\Huge}
\titlespacing*{\chapter}{0pt}{50pt}{40pt}

% Document info
\title{Template Document}
\author{Academic Support Agent}
\date{\today}

\begin{document}
\maketitle

\title{T tulo do Projeto}
  \foreigntitle{Thesis Title}
  \author{Nome}{Sobrenome}
  \advisor{Prof.}{Nome do Primeiro Orientador}{Sobrenome}{D.Sc.}
% \advisor{Prof.}{Nome do Segundo Orientador}{Sobrenome}{D.Sc.}
  
  \examiner{Prof.}{Nome do Primeiro Examinador Sobrenome}{D.Sc.}
  \examiner{Prof.}{Nome do Segundo Examinador Sobrenome}{Ph.D.}
  \examiner{Prof.}{Nome do Terceiro Examinador Sobrenome}{D.Sc.}
  \examiner{Prof.}{Nome do Quarto Examinador Sobrenome}{Dr.-Ing.}
  \examiner{Prof.}{Nome do Quinto Examinador Sobrenome}{M.Sc.}
 
% Para Eng. Mec nica ativar {DEM} e para Eng. Controle e Automa o ativar {DECA} 
  %\department{DEM}
  \department{DECA}
  
  \date{9}{2021}

  \keyword{Primeira palavra-chave}
  \keyword{Segunda palavra-chave}
  \keyword{Terceira palavra-chave}

  \maketitle

  
  % Content from dedic would be included here

  % Content from thanks would be included here

  % Content from resumo would be included here

  % Content from abstract would be included here


  
  
  \printlosymbols
  \printloabbreviations
  
  
  

%  \linenumbers % inserir numera o nas linhas do texto
  % Content from chap01 would be included here

  % Content from chap02 would be included here

  % Content from chap03 would be included here

  % Content from chap04 would be included here

  % Content from chap05 would be included here


  
  
  
  

  
  % Content from appenA would be included here

  % Content from appenB would be included here
\end{document}